\FloatBarrier
\section{Cross-system synthesis: contrasting forecast deficits, one decomposition}\label{sec:synthesis}

The two systems were scored separately and will stay separate; what they share is a question and, it turns out, an arithmetic. Both ladders consist of \emph{increment forecasters}: each module forecasts next year's state as this year's state plus a mechanism, and the mechanism is evaluated at whatever the module can substitute for the quantities it needs but does not know. On cod the stock-flow modules receive catch while the autonomous modules substitute its training mean and retain state-dependent production; on J-17 M2 substitutes last year's recharge, M2m its climatological mean and the conditional diagnostic the realized value. Persistence is the increment forecaster with mechanism zero. The following proposition says exactly how such a forecaster can gain or lose against it.

\subsection{The increment decomposition}\label{sec:decomp}

\begin{proposition}[One-step skill of an increment forecaster]\label{prop:increment}
Let $\Delta_{t+1}=y_{t+1}-y_t$ and let a forecaster issue $\hat y_{t+1}=y_t+m(\hat x_{t+1})$, where $m$ is its mechanism, $x_{t+1}$ the driver at which the mechanism would be evaluated if the driver were known, and $\hat x_{t+1}$ the value substituted at the origin. Write $m_{t+1}=m(x_{t+1})$, $u_{t+1}=\Delta_{t+1}-m_{t+1}$ (the mechanism residual) and $d_{t+1}=m_{t+1}-m(\hat x_{t+1})$ (the substitution error). Then, whenever the second moments exist,
\begin{equation}\label{eq:decomp}
\MSE(\hat y)-\MSE(\pers)\;=\;\underbrace{\E[d_{t+1}^2]-\E[m_{t+1}^2]}_{\text{substitution term}}\;-\;\underbrace{2\,\E\bigl[u_{t+1}\,m(\hat x_{t+1})\bigr]}_{\text{misspecification term}} .
\end{equation}
In particular:
\begin{enumerate}[label=(\roman*),leftmargin=2.2em,itemsep=2pt]
\item \emph{Supplied driver.} If $\hat x_{t+1}=x_{t+1}$ then $d\equiv0$ and $\MSE(\hat y)-\MSE(\pers)=\E[u_{t+1}^2]-\E[\Delta_{t+1}^2]$: the forecaster beats persistence if and only if its mechanism residual is smaller in mean square than the increment itself.
\item \emph{Persisted driver.} If $\hat x_{t+1}=x_t$, the sequence $m_t$ is weakly stationary with mean $\mu$, variance $\sigma^2$ and lag-one autocorrelation $\rho$, and $\E[u_{t+1}m_t]=0$, then
\[\MSE(\hat y)-\MSE(\pers)\;=\;\sigma^2(1-2\rho)-\mu^2,\]
so the forecaster beats persistence if and only if $\rho>\tfrac12-\mu^2/(2\sigma^2)$. For a mean-zero mechanism this is $\rho>\tfrac12$, whatever the mechanism's fidelity.
\item \emph{Constant substituted mechanism.} If the entire forecast increment satisfies $m(\hat x_{t+1})=\mu$ and $\E[u_{t+1}]=0$, then $\MSE(\hat y)-\MSE(\pers)=-\mu^2\le0$: this constant-increment forecaster never loses to persistence, and it beats the persisted-driver forecaster of (ii) by $\sigma^2(1-2\rho)$, which is positive if and only if $\rho<\tfrac12$. A model with climatological flux inputs but a head-dependent drain does not, in general, satisfy this condition.
\end{enumerate}
\end{proposition}

The proof is in \ref{app:proofs}. The identity is exact; its special-case readings require their stated assumptions. When a driver is persisted, the half-autocorrelation threshold follows only for a weakly stationary, mean-zero mechanism with zero residual cross-term. The $-\mu^2$ result applies only when the \emph{entire} substituted increment is constant, not merely when one set of inputs is replaced by its training mean. In general, the substitution term depends on the difference between the realized and substituted mechanisms, while the residual cross-term depends on how the mechanism error co-moves with the forecast actually issued. A structural forecaster can therefore lose to persistence through a substitution penalty partly offset or amplified by residual covariance. Cod supplies the driver at one step, making the misspecification comparison direct; the J-17 rolling application below reports both terms rather than imposing the stationary special case.

\paragraph{J-17: substitution penalty and residual covariance.} For each one-step rolling origin, let $m_{t+1}$ be the fitted M2 increment evaluated at the realized target-year $(R_{t+1},P_{t+1})$ and origin-known $H_t$, and let $\hat m_{t+1}$ use the persisted $(R_t,P_t)$ with the same fitted coefficients and $H_t$. The general identity, evaluated before clipping and with all 75 origin-specific fits, gives an empirical substitution contribution $\overline{d^2-m^2}=+76.03$\,ft$^2$ and residual cross-term $-2\overline{u\hat m}=-27.49$\,ft$^2$, for a raw M2-minus-persistence gap of $+48.54$\,ft$^2$ (SI S23). The $[610,710]$ output clip binds once in this one-step rolling set, at the 1992 origin; the scored gap is $14.70^2-13.23^2=+40.99$\,ft$^2$. Thus driver substitution is the larger adverse term, but it is not the whole forecast-error difference. The full-sample flux component $\hat\beta R_t+\hat\gamma P_t$ has lag-one autocorrelation about 0.08 (standard deviation 9.2\,ft); recharge itself has autocorrelation 0.17, and $r(\Delta H_t,R_t)=0.74$ versus $r(\Delta H_t,R_{t-1})=0.06$. These diagnostics support a low-persistence driver interpretation, but the half-autocorrelation special case is not an exact numerical prediction for this fitted model: its increment also contains the state-dependent drain, and the residual cross-term is nonzero. M2m's 12.28\,ft RMSE is an observed point-score pass; because its affine increment retains the head-dependent drain, it does not instantiate the theorem's constant-increment corollary. The realized-flux score of 7.55\,ft is a conditional diagnostic using future fluxes, not forecast skill. Fig.~\ref{fig:mechanism}a--b shows the separate driver and increment correlations.

\paragraph{Cod fails on the misspecification term.} The stock-flow modules are case (i): the catch is supplied, $d\equiv0$, and their deficit is entirely $\E[u^2]-\E[\Delta^2]$, the mechanism residual against the raw increment. Proposition~\ref{prop:repeller}(ii) says what that residual is in the collapse years: the realised increments $\Delta_{t+1}=S_{t+1}-S_t$ lay 66--161\,kt (coarse regime) and 59--240\,kt (annual landings) below $-C_t$, so $u=\Delta_{t+1}+C_t-g(S_t)<0$ for every nonnegative production curve, and $|u|>|\Delta_{t+1}|$ exactly when $g(S_t)>C_t$, that is, whenever the fitted production exceeds the supplied catch, which is the situation of a module told that removals are falling. The autonomous modules replace catch with its training mean, but their production term remains state-dependent; neither the persisted-driver special case nor the constant-increment corollary applies. After 1992 that substituted catch greatly exceeds production at the depleted stock, the map is driven to the floor, and the residual equals the observed biomass at every origin. Neither failure has anything to do with the predictability of a driver: the apparent net production of the assessed series is strongly autocorrelated (lag-one autocorrelation 0.75 on Specification A under the coarse regime, 0.71 under annual landings, 0.91 on Specification B), and the attribution of the cod result to a ``weakly autocorrelated surplus production'' in an earlier combination of these studies is withdrawn. The cod modules fail because the class cannot represent the transitions that carry the error; the simulations show high sensitivity for the D1 collapse-fit autonomous truth, but low power for the tested stock-flow truth, so non-retention of the catch-supplied modules does not rule out that structure.

\subsection{When climatology overtakes persistence}\label{sec:crossover}

The same arithmetic settles the relation between the two baselines, which the five-year J-17 result raised.

\begin{remark}[Horizon crossover of the baselines]\label{rem:crossover}
If $y$ is weakly stationary with variance $\sigma_y^2$ and lag-$h$ autocorrelation $\rho_h$, and the climatological mean is known, then $\MSE_h(\pers)=2\sigma_y^2(1-\rho_h)$ and $\MSE_h(\mathrm{mean})=\sigma_y^2$, so $\RMSE_h(\pers)/\RMSE_h(\mathrm{mean})=\sqrt{2(1-\rho_h)}$ and the training mean overtakes persistence at the first horizon at which $\rho_h<\tfrac12$. With an estimated mean the ratio is shifted by the estimation variance of the mean, which is small relative to $\sigma_y^2$ for long training windows.
\end{remark}

Fig.~\ref{fig:mechanism}c tests the remark on all three series with the baselines recomputed at every horizon $h=1,\dots,5$ on the registered data and on the origin sets of the scored ladders (the $h=1$ and $h=5$ values reproduce Tables~\ref{tab:rolling} and \ref{tab:edwrolling} exactly; SI S23). On J-17 the head's autocorrelation falls from 0.63 at lag one to 0.35 at lag two, so the remark predicts that climatology overtakes persistence between one and two years; the observed ratios are 0.82, 1.07, 1.22, 1.27 and 1.26 against predicted 0.86, 1.14, 1.31, 1.37 and 1.36, the crossover is at $h=2$, and the five-year result of Section~\ref{sec:edwards} is the plateau of that curve rather than an isolated finding. On cod the SSB series is not stationary over the scored windows (a collapse by an order of magnitude and a partial rebuild), the training mean of a collapsing series is a poor estimate of anything, and the remark fails in the direction it should: the observed ratios stay between 0.18 and 0.58 at every horizon while the stationary formula would predict a crossover between four and five years on Specification A. The criterion is therefore a test as well as a prediction. A close match between $\sqrt{2(1-\hat\rho_h)}$ and the observed baseline ratio indicates that the stationary identity approximates the scored comparison on that series; a mismatch indicates that its assumptions do not describe the series. Neither reading is a policy rule, and neither transfers without a system-specific evaluation.

\begin{figure}[!htb]
\centering
\includegraphics[width=\textwidth]{fig7_mechanism}
\caption{Driver, increment and baseline diagnostics. (a) J-17 annual recharge is weakly persistent: $R_t$ against $R_{t-1}$, $r=0.17$. (b) The J-17 head increment is associated with same-year recharge: $\Delta H_t$ against $R_t$, $r=0.74$; the realized-flux score is conditional on future drivers, and the exact rolling M2 decomposition retains both the substitution and residual-cross terms (SI S23). (c) Ratio of rolling-origin persistence RMSE to training-mean RMSE by horizon on the scored origin sets: J-17 observed (black) against the stationary prediction $\sqrt{2(1-\rho_h)}$ of Remark~\ref{rem:crossover} (grey, dashed), with the crossover at $h=2$; Northern cod Specifications A and B never cross within five years, and the stationary prediction does not apply to them.}
\label{fig:mechanism}
\end{figure}

\subsection{What transfers and what does not}\label{sec:transfer}

Table~\ref{tab:cross} sets the two systems side by side without pooling them. What transfers is an analytic distinction: compare each forecast with origin-matched simple baselines at declared horizons, state the information set and the model-class rule, and keep the exact decomposition's substitution and residual terms separate. On J-17, the realized-flux diagnostic shows how the fitted map scores when supplied target-year drivers, but it is not an operational forecast; the annual-origin climate variants do not recover the recharge pulse in the drought-recovery window. Any within-year data-update forecast would require its own origin definition and evaluation, and no such result is claimed here. On cod, the catch-supplied modules carry a short-memory growth component with one-step skill in a low-biomass band on one specification; this is a hypothesis for a future test with an externally fixed regime boundary and consistent catch timing, not a management recommendation. What does not transfer is any score, margin, autocorrelation or operating characteristic: those results belong to the tested series and information sets, and another system requires its own evaluation.

\begin{table}[!htb]
\centering\small
\caption{The two systems side by side. Nothing in this table is pooled; each column is one series under its own protocol. ``Failing term'' refers to equation~\eqref{eq:decomp}.}
\label{tab:cross}
\begin{tabularx}{\textwidth}{@{}>{\raggedright\arraybackslash}p{0.21\textwidth}LL@{}}
\toprule
 & Northern cod (Spec.\ A; Spec.\ B) & Edwards J-17 \\
\midrule
Predictand & assessed SSB, kt; retrospective vintage & measured annual-mean head, ft \\
Structural claim & surplus production; catch-driven stock flow; depensation; residual memory & one-pool water balance; residual and delay variants; climate-informed recharge \\
Driver at the origin & catch supplied (M2--M4) or held at training mean (M1, M1b) & recharge and pumpage persisted (M2), climatological (M2m), forecast from climate indices, or realised (oracle) \\
Protocol & C: 5\% band, both horizons, comparator and persistence & Pass~1 H1/H2 at $h=1$; separate M2m class disposition; Pass~2 H1/H2 against persistence and M1 \\
One-year persistence RMSE & 98.0; 84.4\,kt & 13.23\,ft \\
Best point-score module at $h=1$ & M1b 114.8 (ratio 1.17); M1 119.5 (1.42) & M2m 12.28 (H1/H2 pass, class-demoted); M1 12.84 (retained) \\
Structural disposition & none; none & Pass~1: M1 retained, M2m class-demoted; Pass~2: three point-gate passes class-demoted; none under C \\
Error decomposition & misspecification: collapse transitions outside the class; floor forecasts after 1992 & M2 substitution term $+76.0$\,ft$^2$, residual cross-term $-27.5$\,ft$^2$; clipped MSE gap $+41.0$\,ft$^2$ \\
Lag-one autocorrelation & apparent net production 0.75; 0.91 & flux component 0.08; recharge 0.17 \\
Baseline crossover (Remark~\ref{rem:crossover}) & none within five years; stationary formula inapplicable & at $h=2$, as predicted on these origins \\
Operating-characteristic evidence & D1 true-M1 power 95.0--97.5\%; D3 true-M2 power 10.5--13.0\%; D5 any-module false retention 2.5--7.5\% (point specificity 92.5--97.5\%) & 7.55\,ft realized-flux conditional diagnostic, not power; climate point margins $\le0.13$\,ft \\
\bottomrule
\end{tabularx}
\end{table}
